\section{Numerical methods, audits, and extended results}
\label{app:numerics}

This section specifies how the protected spaces and tangent responses are
constructed.  Numerical values, including ranks, gaps, residuals, basis
dimensions, and production status, enter through the generated v13 macros and
tables.  The procedures below define those quantities and the gates applied to
them; they do not infer a missing gate from a visually narrow spectrum.

\subsection{Common kernel and response audit}

For a factorized parent, we assemble or apply the rectangular constraint map
$C$ and solve for its right kernel.  For a cohomological parent, we construct
the charge-resolved Hodge Laplacian and solve its zero eigenspace.  For a
stabilizer code, the protected rank follows from binary symplectic elimination
rather than from allocating the exponentially large state space.  In every
dense or sparse eigensolve the accepted projector frame $U_P$ must satisfy
\[
 \lVert U_P^\dagger U_P-I_D\rVert<\epsilon_{\mathrm{orth}},
 \qquad
 \lVert HU_P-E_0U_P\rVert_F<\epsilon_{\mathrm{ker}}.
\]
The observed rank must equal the independently calculated protected rank.  We
then locate the first eigenvalue outside the target fiber and require a
positive gap.  An internal bandwidth smaller than a numerical tolerance is
reported as a residual check; it is not substituted for the algebraic kernel
condition.

Response vectors are obtained by applying the complement inverse to
$Q(\partial_aH)P$.  Dense cases use the stored positive-energy eigenbasis.
Factorized large cases solve shifted or projected linear systems without
forming $H=C^\dagger C$ densely.  Each response is checked for target leakage,
linear-system residual, and stability under the registered solver shift.
Tangent panels are normalized with their own two-point Gram matrix before a
four-channel quantity is evaluated.

\subsection{Kapit--Mueller and continuum $V_0$ parents}

The Kapit--Mueller single-particle hopping is assembled on a square magnetic
torus, including its magnetic-periodic winding phases.  The lowest isolated
band is diagonalized, and its orbital frame is used to map onsite boson pairs
into the projected $N$-particle Fock basis.  The columns of the resulting pair
constraint factor encode the pair annihilators $A_R$ of the $V_0$ parent.
Kernel vectors are computed from the factor,
then checked against the explicitly applied positive parent.  Quasihole ranks
are computed from cyclic $(1,2)$ admissible configurations before the
eigensolve and are compared with the observed nullity.

For the independent continuum check, $N$ bosons occupy $n_\phi$ Landau-gauge
orbitals on a square torus.  Magnetic translates of the periodic Gaussian
give a factorized $V_0$ pair map from the $N$-particle basis to the
$(N-2)$-particle basis.  This code path shares neither the Kapit--Mueller
hopping builder nor its orbital projection.  Agreement of the protected count
therefore checks the two-body clustering mechanism across two microscopic
realizations, whereas disagreement of higher response statistics remains an
allowed parent dependence.

Density tangents are diagonal one-body operators in real space.  Bond and
current tangents are respectively the Hermitian and anti-Hermitian quadratures
of a hopping link.  For the exactness-preserving constraint/generator
transport, the one-body generator is lifted to the fixed-$N$ Fock space and
applied to both the constraint factor and the projector.  We verify $S_a=0$
along the path and compare the transported kernel with a direct kernel solve.
The v6 protected-generator results belong to exact transport.  Only this
calculation supports finite-path and holonomy conclusions.  For twist-bundle
calculations, every mesh edge stores the overlap singular values needed for
polar parallel transport; a mesh is rejected if the target rank or external
gap changes at any twist.

The v2/v3 Laughlin fixed-parent local-potential panels are evaluated
separately.  Their structured or random local potentials are applied only at
the accepted exact parent, and their response is computed from
$Q(\partial_aH)P$ for the isolated rank-$D$ cluster.  These probes generally
do not preserve exact degeneracy away from the base point and therefore do
not establish a finite exact path.  Their admissible conclusions are the
base-point curvature spectrum, covariance, and four-channel statistics.

\subsection{Moore--Read coherent-state constraint map}

In the $N$-boson occupation basis of $N_\phi$ torus orbitals, each coherent
annihilator $b_x=\sum_j\phi_{xj}b_j$ is cubed and divided by $\sqrt{3!}$.
Stacking $C_x$ over an $N_\phi\times N_\phi$ magnetic-translation orbit gives
a sparse map from the $N$-particle Hilbert space to the $(N-3)$-particle
space.  The parent is applied as two sparse multiplications, $v\mapsto
C^\dagger(Cv)$.  The expected rank is computed from cyclic $(2,2)$ admissible
configurations before the solver is run.

The opened cases use periodic boundary conditions and the fixed-four-
quasihole relation $N_\phi=N+2$.  The stored kernel audit contains the basis
dimension, constraint shape and nonzero count, observed and expected ranks,
zero-band width, first positive eigenvalue, residual norm, and orthonormality
error.  The v9 Moore--Read guiding-center panels are exactness-preserving
generator transports.  For each one-body generator $G_a$, the differentiated
constraint-intertwiner identity gives
$(\partial_aH)P=-H\Gamma_N(G_a)P$ and
$X_a=Q\Gamma_N(G_a)P$ in the main-text Kato convention.  The historical v9 routine stores instead
$\widetilde X_a=-Q\Gamma_N(G_a)P=-X_a$, because its solver variable is
$R_Q(\partial_aH)P$.  It checks the parent-action residual of
$\widetilde X_a$, complement leakage, and the condition
\texttt{fiber\_tangent\_zero}, which is the numerical $S_a=0$ test.  The
registered quadratic and fourth-order contractions are invariant under this
common sign, but the distinction is retained for response-factor provenance.
The v9 panels are therefore not generic splitting probes.  Larger staged
cases remain marked pending until the same gates are passed; resource
estimates are not converted into scientific results.

\subsection{Random cubic supercharge}

For $N$ complex fermion modes, a charge-$r$ basis is stored as ordered bit
strings of Hamming weight $r$.  Fermionic parity signs are included when the
cubic annihilation operator
$Q(C)=\sum_{i<j<k}C_{ijk}\psi_i\psi_j\psi_k$ is assembled.
The charge-restricted Hamiltonian is the corresponding block of
$Q^\dagger Q+QQ^\dagger$.  A coupling realization is accepted only if
nilpotency, Hermiticity, the expected cohomology rank, a positive complement,
kernel residual, and orthogonality all pass independently.

Couplings are normalized complex three-forms.  The registered sizes and
central or adjacent charge sectors are stored with independent realization
seeds.  Sparse tangent panels vary a fixed subset of $C_{ijk}$; isotropic
panels vary all cubic couplings with the same norm.  The response is evaluated
both from the direct Hamiltonian derivative and from its exact and coexact
Hodge branches.  We require the two branches to be orthogonal and their sum to
reproduce the direct response.  The random-supercharge calculation has
independent coupling realizations for its registered separable null tests, but
the v13 evidence table leaves complete nonseparable covariance closure
untested.

\subsection{Lattice supercharge on cycle unions}

For $G=\bigsqcup_{s=1}^{m}C_6$, an occupation bit string is retained only when
no graph edge has both endpoints occupied.  The charge-$f$ basis thus
enumerates the size-$f$ faces of the independence complex.  The sparse map
$Q_f:\cH_f\rightarrow\cH_{f+1}$ includes the Jordan--Wigner parity associated
with the canonical vertex order.  Direct multiplication verifies
$Q_{f+1}Q_f=0$.  In sector $f=2m$ we diagonalize
\[
 H_f=Q_{f-1}Q_{f-1}^\dagger+Q_f^\dagger Q_f.
\]
The first $2^m$ eigenvectors must exhaust the numerical kernel, and the next
eigenvalue must be positive.

Complex site couplings are normalized separately on each six-site ring.
Candidate local or isotropic tangent vectors are projected orthogonally to
every component-wise rescaling direction, then orthonormalized.  This removes
the trivial deformation that merely rescales one Hodge block.  Derivatives of
both terms in $H_f$ give exact and coexact response branches.  We compare
their sum with the direct resolvent derivative and check mutual
orthogonality, target leakage, and nonzero branch weights.  The statistical
unit in the complete-covariance calculation is one full fixed-base tangent
panel.  It is not an independent graph or Hamiltonian realization.
For the corrected v13 calculation, the mean of the local-panel population is
obtained by exact enumeration of the ordered coordinate-channel marginals.
The first twelve panels then determine only the channel whitener, while the
remaining twelve panels enter the primary complete-cumulant estimator.

\subsection{X-cube binary-symplectic construction}

Edges of the periodic cubic lattice are indexed by position and one of three
directions.  A Pauli operator is represented by a binary row $(\bm x|\bm z)$.
The symplectic inner product verifies that every cube-$X$ row commutes with
every vertex-plane-$Z$ row.  Gaussian elimination over $\mathbb F_2$ gives the
stabilizer rank, logical-qubit count, and hence the protected rank without
constructing a $2^{3L^3}$-dimensional vector.

For positive coefficient reweighting, the stabilizer table and common $+1$
eigenspace are identical at every parameter value, so $\partial_aP=0$ follows
algebraically.  For the moving-projector control, two local Pauli generators
are selected whose product differs from a stabilizer only by phase.  Their
support, anticommutation, detectability, exact isospectrality, and gap
preservation are checked in the binary representation.  A small dense Pauli
realization independently checks the sign of the curvature.  Production
sizes use the algebraic proof and never allocate the ground-space matrix.

\subsection{Checkerboard low-energy cluster}

The checkerboard control is built from a two-sublattice short-range hopping
matrix on a torus.  Complex nearest-neighbor and real longer-range hoppings
produce an isolated lower Chern band.  At every boundary twist, the lowest
single-particle band frame is recomputed and the onsite boson contact
interaction is projected into its complete fixed-$N$ occupation basis.  This
is the same projected-contact convention used for the growing-rank
Kapit--Mueller comparison, but with an independently constructed hopping
matrix.  Boundary twists are inserted as Peierls phases on seam-crossing bonds
before the projection.  The tracked multiplet is selected by dimension and
continuity over the twist mesh.  Acceptance requires its internal width to
remain smaller than its external separation, but the width is not required to
vanish.

Bond-magnitude, bond-phase, and zero-mean onsite disorder define the
deformation.  Neighboring subspaces are compared by singular values and polar
parallel transport before their curvature increments are formed.  Because the
cluster is only quasi-degenerate, its energy statistics remain meaningful and
are reported separately from the exactly degenerate parents.  No checkerboard
result is used to establish a common-kernel, cohomology, or stabilizer claim.
